\scanpage{132}

précédemment à propos des $N_V,I_V\ldots$), je vais regarder d'un peu plus près l'allure des sous-groupes $I_x\subset\unsure{\pi|\mathcal E}_{L/K}^{\bar L}$, (ou \deleted{\unsure{associés}} plutôt, des classes de \inserted{$\pi$-}conjugaison de tels sous-groupes) associés aux $x\in X\setminus U$, \inserted{pour $x$ variable,} et me \unsure{convaincre} du caractère intrinsèque de ces groupes, indépendamment du choix du modèle $X$ (mais ceci est évident, si on admet ce qu'on vient d'\unsure{énoncer|annoncer} --- que ce sont justement, à sous-groupes ouverts près \deleted{\unsure{tout}} au plus, les images des $I_V$, pour $V\deleted{\unsure{\to}}\inserted{\in}\mathfrak X_{L/K}$ n'ayant \inserted{pas de} centre dans $\deleted{X\setminus U}U\ldots$)

\deleted{Considérons d'abord les}

Soit $J$ une partie de l'ens. fini $I$, et soit
\[
\tag{36}
D_J^{\unsure{\circ}}=\{x\in X\mid I(x)=J\}
=\overbrace{\bigcap_{j\in J}D_j}^{\Delta_J}-\bigcup_{i\in I-J}\Delta_J\cap D_i
\]
\[
D_{\{i\}}=D_i\setminus\left(\bigcup_{j\ne i}\unsure{D_i\cap D_j}\right).
\]
Si $d=\operatorname{card}J$, --- de sorte que l'on voit que $D_J\subset X$ est \inserted{régulier,} de codim $d$ en chacun de ses points (mais bien sûr, il peut être $\varnothing$). Si on appelle ``simplexe'' une partie $J\subset I$

\margin{il vaut mieux définir $D_J$ comme \uline{sous-schéma} \unsure{comme} $\bigcap_{j\in J}D_j$ (36), et définir son \uline{\unsure{support|sous-ouvert}} comme $\{x\in X\mid I(x)\unsure{=|\supseteq}J\}$}

\margin{131}
